% cap1_introducere.tex
\chapter{Introducere}
\label{ch:introducere}

% ─────────────────────────────────────────────────────────────────────────────
\section{Context și motivație}
\label{sec:motivatie}

Locuințele inteligente reprezintă una dintre aplicațiile cu cea mai rapidă
expansiune din domeniul \gls{iot}. Dacă la începutul deceniului trecut
termenul „smart home" evoca un lux rezervat câtorva entuziaști, astăzi
milioane de gospodării integrează senzori de temperatură, camere de
supraveghere, sisteme de iluminat și asistenți vocali conectați la internet.
Piața globală a locuințelor inteligente a depășit \$80 de miliarde în 2023
și se preconizează o creștere susținută în anii următori
\cite{stojkoska2017smarthome,gubbi2013iot}. La baza acestei transformări stă
tocmai convergența \gls{iot}: dispozitive cu resurse limitate care colectează
date din mediul fizic, le transmit în rețea și declanșează acțiuni automatizate
în funcție de starea mediului \cite{atzori2010iot}.

Cu toate acestea, marea majoritate a soluțiilor comerciale disponibile astăzi
(Google Home, Amazon Alexa, Apple HomeKit, Tuya etc.) se bazează pe un model
\emph{cloud-first}: datele brute ale senzorilor sunt trimise pe servere aflate
în altă țară, procesate acolo și returnate aplicației mobile printr-un lanț
de apeluri de rețea. Această arhitectură ridică o serie de probleme grave:

\begin{itemize}
  \item \textbf{Confidențialitate și protecția datelor.}
  Informații sensibile privind comportamentul cotidian al locatarilor
  (ore de prezență acasă, temperaturi, mișcări în cameră) sunt stocate
  pe infrastructuri terțe, adesea în afara jurisdicției \gls{gdpr}
  \cite{zheng2018smarthomeprivacy,gdpr2016}.
  Chiar atunci când furnizorii respectă legislația europeană, utilizatorul
  nu deține controlul efectiv asupra datelor generate în propria locuință
  \cite{alrawi2019sokhome}.

  \item \textbf{Dependența de conectivitatea la internet.}
  Dacă furnizorul de servicii întrerupe suportul pentru un produs sau dacă
  conexiunea la internet este indisponibilă temporar, întregul sistem poate
  deveni nefuncțional. Experiențele legate de retragerea prematură a unor
  platforme (Google Stadia, SmartThings Hub v2 etc.) au demonstrat că
  \emph{vendor lock-in}-ul reprezintă un risc real pentru utilizatorii care
  investesc în ecosisteme proprietare \cite{kleppmann2019localfirst}.

  \item \textbf{Cost și complexitate.}
  Abonamentele lunare pentru platformele cloud, costul gateway-urilor
  proprietare și imposibilitatea de a integra hardware eterogen cresc
  semnificativ bariera de acces. Soluțiile open-source existente (Home
  Assistant, OpenHAB) oferă alternativă, dar cer cunoștințe tehnice
  avansate și nu includ componente de analiză inteligentă a comportamentului.

  \item \textbf{Latență și procesare la margine rețelei.}
  Trimiterea fiecărei citiri de senzor spre cloud și așteptarea unui răspuns
  introduce latențe inacceptabile pentru scenarii de siguranță (detectare gaz,
  alertă de intruziune nocturnă). Paradigma \emph{edge computing} propune
  procesarea locală a datelor, acolo unde sunt generate
  \cite{shi2016edge}.
\end{itemize}

Aceste limitări nu sunt de natură exclusiv tehnică: ele reflectă un dezechilibru
structural în care controlul locuinței proprii este delegat unor entități externe.
Lucrarea de față propune o arhitectură alternativă, \emph{local-first}, care
rezolvă fiecare dintre problemele enumerate anterior.

% ─────────────────────────────────────────────────────────────────────────────
\section{Problema abordată și abordarea local-first}
\label{sec:problema}

Problema centrală abordată în această lucrare este următoarea:
\emph{Este posibil să se construiască un sistem de locuință inteligentă complet
funcțional, capabil de monitorizare în timp real, control al actuatorilor,
supraveghere video și detecție inteligentă a anomaliilor comportamentale,
folosind exclusiv infrastructură locală, fără nicio dependență de servicii cloud,
și cu un cost total de hardware accesibil?}

Răspunsul propus este o arhitectură \textbf{local-first}
\cite{kleppmann2019localfirst}: toate datele rămân în rețeaua locală (\gls{lan})
a locatarului, procesarea se realizează pe un gateway local (\gls{rpi}), iar
aplicația mobilă comunică direct cu serverul intern fără a tranzita internetul.
Această abordare se înscrie în tendința mai largă de \emph{edge computing}
\cite{shi2016edge,premsankar2018edgeiot}: în loc să se centralizeze inteligența
în cloud, logica de decizie este mutată cât mai aproape de sursele de date
(senzori, camere, actuatori).

Sistemul implementat funcționează integral în \gls{lan}, inclusiv atunci când
conexiunea la internet este indisponibilă. Singurul trafic care poate opțional
ieși din rețeaua locală este notificarea push, iar și aceasta poate fi dezactivată.
Datele personale ale utilizatorului (temperaturi, ore de prezență, imagini de la
cameră) nu părăsesc niciodată locuința fără consimțământ explicit.

Componenta de \gls{ml} — motorul de detecție a anomaliilor comportamentale bazat
pe algoritmul \gls{if} — rulează direct pe \gls{rpi}, eliminând necesitatea
unui serviciu cloud de inferență. Aceasta este contribuția originală centrală a
lucrării, detaliată în \secref{sec:contributii} și dezvoltată extensiv în
capitolele \chapref{ch:metodologie} și \chapref{ch:implementare}.

% ─────────────────────────────────────────────────────────────────────────────
\section{Obiectivele lucrării}
\label{sec:obiective}

Lucrarea își propune proiectarea, implementarea și evaluarea unui sistem de
locuință inteligentă local-first, cu scopuri concrete structurate pe cinci
straturi arhitecturale. Obiectivele principale sunt:

\begin{enumerate}
  \item \textbf{Sistem funcțional pe cinci straturi.}
  Proiectarea unei arhitecturi coerente care integrează: senzori și actuatori
  fizici, noduri \gls{iot} bazate pe microcontrolere ESP32 programate în
  Arduino C++, un gateway local pe \gls{rpi}, un \gls{api} \gls{rest} +
  WebSocket și o aplicație mobilă — toate comunicând în \gls{lan} fără cloud.

  \item \textbf{Achiziție de date cu hardware real.}
  Conectarea și programarea celor 11 componente hardware reale ale nodului
  interior (ESP32 WROOM-32 cu două senzoare \gls{dht}11, un senzor de gaz
  MQ-2, două rezistențe fotosensibile \gls{ldr}, trei \gls{led}-uri, un
  buzzer, un servomotor SG90 și un ventilator 5V) și a nodului exterior
  (ESP32-CAM cu senzor \gls{pir} și cameră OV2640).

  \item \textbf{Control al actuatorilor în timp real.}
  Implementarea unui protocol de comandă bidirecțional prin \gls{mqtt} care
  permite aprinderea/stingerea \gls{led}-urilor, pornirea ventilatorului,
  mișcarea servo-ului (draperii), activarea buzzerului și configurarea \gls{wifi}
  — toate inițiate din aplicația mobilă cu latență sub o secundă în \gls{lan}.

  \item \textbf{Motor de \gls{ml} local pentru detecția anomaliilor.}
  Implementarea unui pipeline complet de \gls{if} (antrenare, normalizare,
  inferență în timp real, reantrenare automată zilnică) care rulează exclusiv
  pe \gls{rpi}, fără apeluri externe, și care generează alerte de tip
  \texttt{ML\_ANOMALY} integrate în sistemul de notificări al aplicației mobile.

  \item \textbf{Aplicație mobilă cross-platform.}
  Dezvoltarea unei aplicații React Native / Expo cu patru tab-uri funcționale
  (Dashboard cu date live, Istoric cu grafice, Alerte, Control actuatori +
  cameră), care se conectează automat la gateway prin \gls{mdns} sau adresă IP
  și oferă autentificare \gls{jwt}.

  \item \textbf{Independență completă de cloud.}
  Demonstrarea că întregul sistem (monitorizare, control, \gls{ml}, video)
  funcționează fără nicio dependență de servicii externe, asigurând
  confidențialitatea datelor și reziliența la lipsa conectivității la internet.

  \item \textbf{Cost redus și reproductibilitate.}
  Construirea sistemului din componente disponibile comercial, cu un cost total
  de hardware accesibil, și documentarea completă a firmware-ului, configurației
  serverului și structurii bazei de date pentru a permite reproducerea independentă.
\end{enumerate}

% ─────────────────────────────────────────────────────────────────────────────
\section{Contribuții originale}
\label{sec:contributii}

Principala contribuție originală a lucrării este \textbf{implementarea unui motor
de detecție a anomaliilor comportamentale bazat pe algoritmul
\gls{if} \cite{liu2008isolation}, rulat în întregime local pe \gls{rpi}}.

Algoritmul \gls{if}, propus de Liu, Ting și Zhou \cite{liu2008isolation}, izolează
punctele anomale printr-un mecanism de partiționare aleatoare a spațiului de
caracteristici: punctele rare (anomale) sunt separate prin mai puțini pași
față de punctele normale, rezultând o \emph{lungime de cale medie} mai mică și,
implicit, un scor de anomalie mai mare. Această proprietate îl face deosebit de
potrivit pentru scenarii de tip \emph{edge}: nu necesită etichete de antrenare
(algoritm nesupervizat), are complexitate computațională redusă și poate fi
reantrenant periodic pe date noi fără intervenție umană.

În contextul acestui sistem, motorul \gls{ml} lucrează cu 12 trăsături extrase
din citirile senzorilor (temperaturi de la ambele senzoare DHT11, umiditate,
nivel de gaz normalizat, intensitate luminoasă normalizată logaritmic, prezența
mișcării, ora zilei, ziua săptămânii, indicator de weekend/noapte, variația
temperaturii pe o oră și numărul de evenimente de mișcare în 10 minute).
Modelul este antrenat pe ultimele 30 de zile de date și reantrenant automat
zilnic la ora 03:00 printr-un serviciu systemd. Inferența se realizează în timp
real la fiecare citire nouă de senzori, fără latența unui apel de rețea spre
cloud.

Această abordare reprezintă un avans față de soluțiile comerciale existente, care
fie nu includ deloc componente de \gls{ml}, fie externalizează procesarea spre
servicii cloud cu implicații de confidențialitate. Spre deosebire de metodele
supervizate care necesită date etichetate (greu de obținut în medii rezidențiale),
\gls{if} este \emph{nesupervizat} și se adaptează automat la rutina specifică a
locatarilor \cite{chandola2009anomaly,fahim2019anomaly,cook2020anomalyiot}.

Contribuțiile secundare ale lucrării includ:
\begin{itemize}
  \item \textbf{Integrarea completă firmware↔broker↔API↔aplicație mobilă.}
  Lanțul complet de la firmware Arduino C++ pe ESP32, prin brokerul Mosquitto
  \gls{mqtt}, prin \gls{api}-ul NestJS cu TypeORM și TimescaleDB, până la
  interfața React Native — implementat și funcțional pe hardware fizic.

  \item \textbf{Arhitectura de date time-series.}
  Utilizarea TimescaleDB (extensie PostgreSQL orientată pe serii de timp) pentru
  stocarea eficientă a citirilor de senzori, cu indexare automată pe timestamp
  și suport pentru interogări de agregare pe intervale temporale.

  \item \textbf{Provisioning Wi-Fi și descoperire automată în rețea.}
  Mecanism de configurare inițială prin hotspot propriu al nodului ESP32
  (\texttt{SmartHome-Setup}), cu stocare a credențialelor în \gls{nvs} și
  descoperire automată a gateway-ului prin \gls{mdns} (\texttt{smarthome.local}).

  \item \textbf{Modul de simulare pentru dezvoltare fără hardware.}
  Un modul \texttt{simulator/} integrat în backend-ul NestJS care generează date
  realiste ale senzorilor (inclusiv variabilele extrase ulterior de motorul \gls{ml}),
  permițând dezvoltarea și testarea completă a stivei software fără acces la
  dispozitivele fizice.
\end{itemize}

% ─────────────────────────────────────────────────────────────────────────────
\section{Structura lucrării}
\label{sec:structura}

Lucrarea este organizată în șase capitole, fiecare abordând un aspect distinct
al sistemului:

\textbf{\chapref{ch:introducere} — Introducere} (capitolul de față) prezintă
contextul domeniului, motivațiile care au condus la abordarea local-first,
obiectivele concrete ale sistemului și contribuțiile originale.

\textbf{\chapref{ch:studii} — Studiu bibliografic} trece în revistă literatura
relevantă: evoluția sistemelor \gls{iot} și a locuințelor inteligente, protocoale
de comunicație cu accent pe \gls{mqtt}, tehnici de detecție a anomaliilor în date
de senzori (comparație între \gls{if}, LOF, One-Class SVM și metode bazate pe
rețele neuronale), paradigma \emph{edge computing} și cadrul legal \gls{gdpr}
privind datele personale în medii rezidențiale.

\textbf{\chapref{ch:metodologie} — Metodologie și proiectare} descrie arhitectura
pe cinci straturi a sistemului, deciziile de proiectare (alegerea componentelor
hardware, a stivei software, a bazei de date time-series), modelarea entităților,
fluxurile de date \gls{mqtt} și pipeline-ul complet al motorului \gls{ml}.

\textbf{\chapref{ch:implementare} — Implementare} detaliază realizarea efectivă a
fiecărui strat: firmware-ul Arduino C++ al nodului interior și al ESP32-CAM,
configurarea brokerului Mosquitto și a \gls{rpi} ca gateway, implementarea
\gls{api}-ului NestJS cu modulele de \gls{mqtt}, WebSocket, autentificare \gls{jwt}
și \gls{ml}, și dezvoltarea aplicației mobile React Native cu toate funcționalitățile
sale.

\textbf{\chapref{ch:rezultate} — Rezultate și evaluare} prezintă rezultatele
obținute: funcționarea sistemului pe hardware real, evaluarea motorului de \gls{ml}
pe date ilustrative (cu date preliminare marcate explicit ca atare), latențele
măsurate în \gls{lan} și o discuție critică a limitărilor actuale și a direcțiilor
de îmbunătățire.

\textbf{\chapref{ch:concluzii} — Concluzii} sintetizează contribuțiile lucrării,
evaluează măsura în care obiectivele au fost atinse și propune direcții de
cercetare și dezvoltare ulterioară, inclusiv extinderea la mai multe noduri,
integrarea unor senzori suplimentari și posibilitatea unui model \gls{ml} mai
sofisticat (autoencodere, LSTM) pentru detecția anomaliilor.
